% 海河少女系统架构图
% 接口：主文已定义 Ink、JinmenBlue、BrickGrey、Cinnabar、PaleJade、RicePaper。
% 在 figure 环境中直接使用 % 海河少女系统架构图
% 接口：主文已定义 Ink、JinmenBlue、BrickGrey、Cinnabar、PaleJade、RicePaper。
% 在 figure 环境中直接使用 % 海河少女系统架构图
% 接口：主文已定义 Ink、JinmenBlue、BrickGrey、Cinnabar、PaleJade、RicePaper。
% 在 figure 环境中直接使用 % 海河少女系统架构图
% 接口：主文已定义 Ink、JinmenBlue、BrickGrey、Cinnabar、PaleJade、RicePaper。
% 在 figure 环境中直接使用 \input{figures/system-architecture}。
\begin{tikzpicture}[x=1mm,y=1mm]
  \path[use as bounding box] (0,0) rectangle (158,60);

  \tikzset{
    hhArchBase/.style={
      draw=BrickGrey,
      fill=white,
      line width=.55pt,
      rounded corners=1.2mm,
      minimum width=31mm,
      minimum height=14mm,
      text width=27mm,
      inner sep=1.5mm,
      align=center,
      font=\small\sffamily,
      text=Ink
    },
    hhArchCore/.style={
      hhArchBase,
      draw=JinmenBlue,
      line width=.8pt,
      fill=PaleJade!62
    },
    hhArchGuard/.style={
      hhArchBase,
      draw=Cinnabar,
      line width=.8pt,
      fill=RicePaper
    },
    hhArchFlow/.style={
      -{Latex[length=2.1mm,width=1.45mm]},
      draw=JinmenBlue,
      line width=.75pt,
      line cap=round,
      line join=round
    },
    hhArchReject/.style={
      -{Latex[length=2.1mm,width=1.45mm]},
      draw=Cinnabar,
      line width=.75pt,
      line cap=round,
      line join=round
    }
  }

  % 上层：从受控来源到带引用的结构化回答。
  \node[hhArchBase]  (source) at (16,50)
    {\textbf{审核来源}\\[-.2ex]\scriptsize PDF\,/\,DOCX\,/\,MD};
  \node[hhArchCore]  (index)  at (57,50)
    {\textbf{混合索引}\\[-.2ex]\footnotesize FTS5 + 向量};
  \node[hhArchGuard] (gate)   at (98,50)
    {\textbf{证据门控}\\[-.2ex]\footnotesize 范围・阈值・来源};
  \node[hhArchCore]  (gen)    at (139,50)
    {\textbf{结构化生成}\\[-.2ex]\footnotesize 事实段绑定引用};

  % 下层反向回流：复核、授权、流式传输与数字人呈现。
  \node[hhArchBase] (front)  at (16,8)
    {\textbf{文化向导}\\[-.2ex]\footnotesize React + VRM};
  \node[hhArchCore] (sse)    at (57,8)
    {\textbf{流式传输}\\[-.2ex]\footnotesize 字幕・引用・状态};
  \node[hhArchCore] (token)  at (98,8)
    {\textbf{播报授权}\\[-.2ex]\footnotesize 一次性语音令牌};
  \node[hhArchGuard] (verify) at (139,8)
    {\textbf{回答复核}\\[-.2ex]\footnotesize 引用・数字・语义};

  \draw[hhArchFlow] (source.east) -- (index.west);
  \draw[hhArchFlow] (index.east)  -- (gate.west);
  \draw[hhArchFlow] (gate.east)   -- (gen.west);
  \draw[hhArchFlow] (gen.south)   -- (verify.north);
  \draw[hhArchFlow] (verify.west) -- (token.east);
  \draw[hhArchFlow] (token.west)  -- (sse.east);
  \draw[hhArchFlow] (sse.west)    -- (front.east);

  % 证据不足时不进入生成与语音授权，但仍经 SSE 向前端返回可见状态。
  \node[hhArchGuard,minimum width=27mm,minimum height=11mm,text width=23mm]
    (reject) at (98,29) {\textbf{固定拒答}\\[-.2ex]\footnotesize 不进入 TTS};
  \draw[hhArchReject] (gate.south) -- (reject.north);
  \draw[hhArchReject] (reject.west) -- ++(-8mm,0) -| (sse.north);
\end{tikzpicture}
。
\begin{tikzpicture}[x=1mm,y=1mm]
  \path[use as bounding box] (0,0) rectangle (158,60);

  \tikzset{
    hhArchBase/.style={
      draw=BrickGrey,
      fill=white,
      line width=.55pt,
      rounded corners=1.2mm,
      minimum width=31mm,
      minimum height=14mm,
      text width=27mm,
      inner sep=1.5mm,
      align=center,
      font=\small\sffamily,
      text=Ink
    },
    hhArchCore/.style={
      hhArchBase,
      draw=JinmenBlue,
      line width=.8pt,
      fill=PaleJade!62
    },
    hhArchGuard/.style={
      hhArchBase,
      draw=Cinnabar,
      line width=.8pt,
      fill=RicePaper
    },
    hhArchFlow/.style={
      -{Latex[length=2.1mm,width=1.45mm]},
      draw=JinmenBlue,
      line width=.75pt,
      line cap=round,
      line join=round
    },
    hhArchReject/.style={
      -{Latex[length=2.1mm,width=1.45mm]},
      draw=Cinnabar,
      line width=.75pt,
      line cap=round,
      line join=round
    }
  }

  % 上层：从受控来源到带引用的结构化回答。
  \node[hhArchBase]  (source) at (16,50)
    {\textbf{审核来源}\\[-.2ex]\scriptsize PDF\,/\,DOCX\,/\,MD};
  \node[hhArchCore]  (index)  at (57,50)
    {\textbf{混合索引}\\[-.2ex]\footnotesize FTS5 + 向量};
  \node[hhArchGuard] (gate)   at (98,50)
    {\textbf{证据门控}\\[-.2ex]\footnotesize 范围・阈值・来源};
  \node[hhArchCore]  (gen)    at (139,50)
    {\textbf{结构化生成}\\[-.2ex]\footnotesize 事实段绑定引用};

  % 下层反向回流：复核、授权、流式传输与数字人呈现。
  \node[hhArchBase] (front)  at (16,8)
    {\textbf{文化向导}\\[-.2ex]\footnotesize React + VRM};
  \node[hhArchCore] (sse)    at (57,8)
    {\textbf{流式传输}\\[-.2ex]\footnotesize 字幕・引用・状态};
  \node[hhArchCore] (token)  at (98,8)
    {\textbf{播报授权}\\[-.2ex]\footnotesize 一次性语音令牌};
  \node[hhArchGuard] (verify) at (139,8)
    {\textbf{回答复核}\\[-.2ex]\footnotesize 引用・数字・语义};

  \draw[hhArchFlow] (source.east) -- (index.west);
  \draw[hhArchFlow] (index.east)  -- (gate.west);
  \draw[hhArchFlow] (gate.east)   -- (gen.west);
  \draw[hhArchFlow] (gen.south)   -- (verify.north);
  \draw[hhArchFlow] (verify.west) -- (token.east);
  \draw[hhArchFlow] (token.west)  -- (sse.east);
  \draw[hhArchFlow] (sse.west)    -- (front.east);

  % 证据不足时不进入生成与语音授权，但仍经 SSE 向前端返回可见状态。
  \node[hhArchGuard,minimum width=27mm,minimum height=11mm,text width=23mm]
    (reject) at (98,29) {\textbf{固定拒答}\\[-.2ex]\footnotesize 不进入 TTS};
  \draw[hhArchReject] (gate.south) -- (reject.north);
  \draw[hhArchReject] (reject.west) -- ++(-8mm,0) -| (sse.north);
\end{tikzpicture}
。
\begin{tikzpicture}[x=1mm,y=1mm]
  \path[use as bounding box] (0,0) rectangle (158,60);

  \tikzset{
    hhArchBase/.style={
      draw=BrickGrey,
      fill=white,
      line width=.55pt,
      rounded corners=1.2mm,
      minimum width=31mm,
      minimum height=14mm,
      text width=27mm,
      inner sep=1.5mm,
      align=center,
      font=\small\sffamily,
      text=Ink
    },
    hhArchCore/.style={
      hhArchBase,
      draw=JinmenBlue,
      line width=.8pt,
      fill=PaleJade!62
    },
    hhArchGuard/.style={
      hhArchBase,
      draw=Cinnabar,
      line width=.8pt,
      fill=RicePaper
    },
    hhArchFlow/.style={
      -{Latex[length=2.1mm,width=1.45mm]},
      draw=JinmenBlue,
      line width=.75pt,
      line cap=round,
      line join=round
    },
    hhArchReject/.style={
      -{Latex[length=2.1mm,width=1.45mm]},
      draw=Cinnabar,
      line width=.75pt,
      line cap=round,
      line join=round
    }
  }

  % 上层：从受控来源到带引用的结构化回答。
  \node[hhArchBase]  (source) at (16,50)
    {\textbf{审核来源}\\[-.2ex]\scriptsize PDF\,/\,DOCX\,/\,MD};
  \node[hhArchCore]  (index)  at (57,50)
    {\textbf{混合索引}\\[-.2ex]\footnotesize FTS5 + 向量};
  \node[hhArchGuard] (gate)   at (98,50)
    {\textbf{证据门控}\\[-.2ex]\footnotesize 范围・阈值・来源};
  \node[hhArchCore]  (gen)    at (139,50)
    {\textbf{结构化生成}\\[-.2ex]\footnotesize 事实段绑定引用};

  % 下层反向回流：复核、授权、流式传输与数字人呈现。
  \node[hhArchBase] (front)  at (16,8)
    {\textbf{文化向导}\\[-.2ex]\footnotesize React + VRM};
  \node[hhArchCore] (sse)    at (57,8)
    {\textbf{流式传输}\\[-.2ex]\footnotesize 字幕・引用・状态};
  \node[hhArchCore] (token)  at (98,8)
    {\textbf{播报授权}\\[-.2ex]\footnotesize 一次性语音令牌};
  \node[hhArchGuard] (verify) at (139,8)
    {\textbf{回答复核}\\[-.2ex]\footnotesize 引用・数字・语义};

  \draw[hhArchFlow] (source.east) -- (index.west);
  \draw[hhArchFlow] (index.east)  -- (gate.west);
  \draw[hhArchFlow] (gate.east)   -- (gen.west);
  \draw[hhArchFlow] (gen.south)   -- (verify.north);
  \draw[hhArchFlow] (verify.west) -- (token.east);
  \draw[hhArchFlow] (token.west)  -- (sse.east);
  \draw[hhArchFlow] (sse.west)    -- (front.east);

  % 证据不足时不进入生成与语音授权，但仍经 SSE 向前端返回可见状态。
  \node[hhArchGuard,minimum width=27mm,minimum height=11mm,text width=23mm]
    (reject) at (98,29) {\textbf{固定拒答}\\[-.2ex]\footnotesize 不进入 TTS};
  \draw[hhArchReject] (gate.south) -- (reject.north);
  \draw[hhArchReject] (reject.west) -- ++(-8mm,0) -| (sse.north);
\end{tikzpicture}
。
\begin{tikzpicture}[x=1mm,y=1mm]
  \path[use as bounding box] (0,0) rectangle (158,60);

  \tikzset{
    hhArchBase/.style={
      draw=BrickGrey,
      fill=white,
      line width=.55pt,
      rounded corners=1.2mm,
      minimum width=31mm,
      minimum height=14mm,
      text width=27mm,
      inner sep=1.5mm,
      align=center,
      font=\small\sffamily,
      text=Ink
    },
    hhArchCore/.style={
      hhArchBase,
      draw=JinmenBlue,
      line width=.8pt,
      fill=PaleJade!62
    },
    hhArchGuard/.style={
      hhArchBase,
      draw=Cinnabar,
      line width=.8pt,
      fill=RicePaper
    },
    hhArchFlow/.style={
      -{Latex[length=2.1mm,width=1.45mm]},
      draw=JinmenBlue,
      line width=.75pt,
      line cap=round,
      line join=round
    },
    hhArchReject/.style={
      -{Latex[length=2.1mm,width=1.45mm]},
      draw=Cinnabar,
      line width=.75pt,
      line cap=round,
      line join=round
    }
  }

  % 上层：从受控来源到带引用的结构化回答。
  \node[hhArchBase]  (source) at (16,50)
    {\textbf{审核来源}\\[-.2ex]\scriptsize PDF\,/\,DOCX\,/\,MD};
  \node[hhArchCore]  (index)  at (57,50)
    {\textbf{混合索引}\\[-.2ex]\footnotesize FTS5 + 向量};
  \node[hhArchGuard] (gate)   at (98,50)
    {\textbf{证据门控}\\[-.2ex]\footnotesize 范围・阈值・来源};
  \node[hhArchCore]  (gen)    at (139,50)
    {\textbf{结构化生成}\\[-.2ex]\footnotesize 事实段绑定引用};

  % 下层反向回流：复核、授权、流式传输与数字人呈现。
  \node[hhArchBase] (front)  at (16,8)
    {\textbf{文化向导}\\[-.2ex]\footnotesize React + VRM};
  \node[hhArchCore] (sse)    at (57,8)
    {\textbf{流式传输}\\[-.2ex]\footnotesize 字幕・引用・状态};
  \node[hhArchCore] (token)  at (98,8)
    {\textbf{播报授权}\\[-.2ex]\footnotesize 一次性语音令牌};
  \node[hhArchGuard] (verify) at (139,8)
    {\textbf{回答复核}\\[-.2ex]\footnotesize 引用・数字・语义};

  \draw[hhArchFlow] (source.east) -- (index.west);
  \draw[hhArchFlow] (index.east)  -- (gate.west);
  \draw[hhArchFlow] (gate.east)   -- (gen.west);
  \draw[hhArchFlow] (gen.south)   -- (verify.north);
  \draw[hhArchFlow] (verify.west) -- (token.east);
  \draw[hhArchFlow] (token.west)  -- (sse.east);
  \draw[hhArchFlow] (sse.west)    -- (front.east);

  % 证据不足时不进入生成与语音授权，但仍经 SSE 向前端返回可见状态。
  \node[hhArchGuard,minimum width=27mm,minimum height=11mm,text width=23mm]
    (reject) at (98,29) {\textbf{固定拒答}\\[-.2ex]\footnotesize 不进入 TTS};
  \draw[hhArchReject] (gate.south) -- (reject.north);
  \draw[hhArchReject] (reject.west) -- ++(-8mm,0) -| (sse.north);
\end{tikzpicture}
